Reliable Taylor-cone angle identification requires both an accurate field solution and an observable that distinguishes neighbouring candidate geometries. This study compares amplitude-projected stress-residual minimization with flank-field power-law exponent crossing in a reduced axisymmetric electrostatic--capillary model. Sparse finite differences with an immersed interface provide the numerical field solution; manufactured-solution and curved-boundary refinement tests assess parts of that workflow. The residual objective is shallow near its minimum, so distinct cone geometries can give similar values. By contrast, when the analytical Taylor potential is imposed on the far boundary, the candidate angle at which the fitted flank exponent crosses $p=-1$ lies between $49.07^\circ$ and $49.25^\circ$ across the tested settings. This is an ideal-limit consistency result, not an independent prediction of a physical cone angle: with a grounded outer boundary, the fitted exponent never reaches $-1$ and the crossing diagnostic is undefined. A synthetic-image round trip qualifies the extraction software on prescribed analytic silhouettes; a separate workflow applies image extraction and a fixed-pressure stress diagnosis to published cone-jet photographs. The latter is a secondary-data application, not validation of exponent crossing or a calibrated physical measurement. The result is a bounded comparison of angle observables and their boundary-condition dependence.